\subsection{Extracting the rank from a lacunary Hodge type}

\begin{proposition}[Rank extraction in every component]\label{famgal:rank}
Let \(\Gamma\) be profinite and \(V\) a continuous semisimple representation
over \(k=\overline{\mathbf Q}_\ell\). Suppose that on a dense subset the
operator on \(V\) is annihilated by a monic polynomial of degree \(a\ge1\).
Suppose that the algebraic monodromy group of \(V\), after scalar extension
to \(\mathbf C_\ell\), contains a cocharacter for which \(a\) distinct
weights occur. Assume that no two disjoint nonempty subsets of these
\(a\) integers have the same arithmetic mean. There is a semisimple
\(k\)-representation \(W\) of dimension \(a\) with
\(W\hookrightarrow V\hookrightarrow W^{\oplus h}\) for some \(h\).
The elements of \(\Gamma\) having simple spectrum on \(W\) form a dense
open subset.
\end{proposition}
\begin{proof}
Write \(G\) for the Zariski closure of the image and \(H=G^\circ\).
Every element \(A\) of \(G\) has minimal polynomial of degree at most
\(a\). Indeed, linear dependence of \(1,A,\ldots,A^a\) is the
vanishing of minors, and it holds first on the given dense set.
These equations persist over \(\mathbf C_\ell\). The representation of
\(G\) is semisimple since its invariant subspaces are exactly those of
\(\Gamma\). The unipotent radical of \(H\) is normal in \(G\), and its
fixed space in every irreducible \(G\)-summand is nonzero by triangularization.
It thus acts trivially on the faithful representation \(V\), so \(H\)
is reductive.

Choose a maximal torus \(T\) and a Borel subgroup containing it over \(k\).
By conjugacy of maximal tori, we may evaluate the given cocharacter on
the \(T\)-weights after extending scalars. At a general point of \(T\)
all distinct weights give distinct eigenvalues, since their pairwise
quotients are nontrivial characters. There are at most \(a\) weights
and the cocharacter separates at least \(a\). Thus there are exactly
\(a\), all separated, and their cocharacter values are the given set.

In an irreducible \(H\)-module all weights differ by elements of the
root lattice. Averaging over the Weyl group annihilates each root,
since its reflection negates it. The average of each Weyl orbit in
this module is consequently the same rational character. Two different
orbits would give disjoint nonempty subsets of the cocharacter weights
with equal means. There is therefore just one orbit. Its highest weight
has multiplicity one, and Weyl invariance of multiplicities shows that
every weight in the orbit has multiplicity one. If two irreducibles
share a weight they share the orbit and its dominant highest weight,
so they are isomorphic. We use here the ordinary highest-weight theorem
for connected reductive groups in characteristic zero.

We check the action of the components on the multiplicity spaces.
Let \(n\in G(k)\) normalize \(T\) and the chosen Borel. It permutes
the \(a\) weights. For a cycle \(C\) of length \(d_C\), put
\(w_C=\prod_{\alpha\in C}\alpha\). On one weight space \(V_\alpha\)
in the cycle, \((tn)^{d_C}\) acts as \(w_C(t)n^{d_C}\), \(t\in T(k)\).
The spectrum of \(tn\) on the cycle contains all \(d_C\)-th roots of
\(w_C(t)c\) for each eigenvalue \(c\) of \(n^{d_C}|_{V_\alpha}\).
Indeed, propagate an eigenvector around the cycle, dividing by the
chosen root at each step. The resulting components close up on return.
Two different cycles \(C,D\) have disjoint spectra for general \(t\).
An overlap would impose
\(w_C(t)^{d_D}/w_D(t)^{d_C}=c_D^{d_C}/c_C^{d_D}\) for a pair of eigenvalues.
The character on the left is nontrivial by the distinct arithmetic
means on the cocharacter. There are only finitely many pairs.
Thus \(tn\) has at least \(\sum_C d_C=a\) distinct eigenvalues.
Its minimal polynomial has degree at most \(a\), so it is diagonalizable
with exactly this many eigenvalues. Each \(n^{d_C}|_{V_\alpha}\)
is diagonalizable and has only one eigenvalue, hence is scalar.
In particular \(n\) acts by a scalar on any weight space that it fixes.

In each orbit of \(G/H\) on the \(H\)-isotypic spaces choose one
representative and let \(\alpha\) be its highest weight.
The stabilizer of this isotypic space preserves the \(H\)-submodule
generated by any vector of \(V_\alpha\). Indeed, write a stabilizing
element as \(h n\) with \(h\in H\) and \(n\) normalizing the Borel
and torus, using conjugacy of Borel--torus pairs in \(H\).
The element \(n\) fixes \(\alpha\) and acts by a common scalar on
\(V_\alpha\). An \(H\)-intertwiner taking one highest vector to another
also commutes with the stabilizer: on \(h'v\) this follows by writing
\(n h'v=(n h'n^{-1})nv\). A basis of \(V_\alpha\) therefore gives
isomorphic stabilizer modules whose sum is the chosen isotypic space.
Transport each of these modules to the other isotypic spaces in its
orbit. Choice of coset representatives has no effect on the transported
subspaces, and the same intertwiners transport consistently.
Their orbit sums give isomorphic \(G\)-modules. In each orbit retain
one such sum and take their direct sum for \(W\). Every \(T\)-weight
now occurs exactly once. This gives \(\dim W=a\) and the two
injections, taking \(h\) to be the maximum multiplicity.
As a summand of \(V\), \(W\) is semisimple and continuous.

Every component has a representative \(n\) normalizing the Borel and
torus. The cycle calculation on its one-dimensional weight spaces
gives simple spectrum on \(W\) for general \(tn\). The nonvanishing
of the discriminant is thus a nonempty Zariski open set in every
component. The preimage of \(H\) in \(\Gamma\) is open and closed.
Its image is Zariski dense in \(H\), as is the image of each finite-index
subgroup of it, since an irreducible variety cannot be a finite union
of proper closed subsets. Each open subset of \(\Gamma\) contains a
coset of such an open subgroup. Its image meets the discriminant open
in its component. Continuity gives openness on \(\Gamma\) as well.
\end{proof}

For the seed above, use the labelled Hodge--Tate decomposition of the
geometric witness at its distinguished embedding. Its grading defines
a cocharacter in the local algebraic monodromy group, hence in \(G\).
To see the membership, the Hodge--Tate decomposition respects tensor
products, duals, morphisms and subrepresentations of Hodge--Tate
representations. The grading therefore preserves each stable line in
tensor constructions, including direct sums of such constructions.
The stabilizer description of an algebraic subgroup of a general linear
group (Chevalley's theorem) puts the grading in the scalar extension of
that subgroup. With finite coefficients, take the component of
\(\mathbf C_\ell\otimes_{\mathbf Q_\ell}L\) at the embedding giving the
specified label, extending the local embedding to \(\mathbf C_\ell\)
accordingly. This is the labelled cocharacter construction of
Dospinescu, Pa\v{s}k\=unas and Schraen~\cite{senCharacters}.
Changing the sign convention for all weights preserves lacunarity.
Proposition~\ref{famgal:rank} now supplies \eqref{newgal:injections}
for every such witness. The determinant identification is then the
dense Frobenius calculation in Proposition~\ref{newgal:seed}.
The construction still needs the geometric witness and the simultaneous
Hecke congruences at the claimed weights and places.
\src{src/Development/AutomorphicArithmetic/Foundations/GaloisRepresentations/IntegralStructures/ZetaCharacterDivisibility/HtCocharOf.lean}{20}{Hodge cocharacter in the closure},
\src{src/Development/EtaleCohomology/DivisorCycles/GoalHolds.lean}{20}{algebraic rank-extraction input}, and
\src{src/Development/AutomorphicArithmetic/Foundations/GaloisRepresentations/IntegralStructures/ZetaCharacterDivisibility/PDivSet.lean}{16}{application to the Galois seed}.
